\documentclass[10pt,a4paper]{amsart}
\usepackage[margin=19mm]{geometry}
\usepackage{amsmath,amssymb,mathtools}
\usepackage[scr=rsfs,cal=euler]{mathalfa}
\usepackage{xfrac}
\usepackage[unicode,hidelinks,bookmarks=false]{hyperref}
\newcommand{\ie}{i.e.\ }
\newcommand{\cf}{cf.\ }
\newcommand{\resp}{respectively\ }
\newcommand{\Subsection}{\subsection}
\providecommand{\nobreakdash}{\nobreak\hbox{-}}
\setlength{\emergencystretch}{2em}
\multlinegap0pt
\usepackage{amssymb}
\usepackage{tcolorbox}
\usepackage[shortlabels]{enumitem}
\usepackage{float}
\usepackage{xcolor}
\newtheorem{theorem}{Theorem}
\title{On the planar Taylor-Couette system and related exterior problems}
\author{Filippo GAZZOLA -- Ji\v{r}\'{\i} NEUSTUPA -- Gianmarco SPERONE}
\date{Source: arXiv:2406.14960v2 (CC BY 4.0)}
\begin{document}
\maketitle

\noindent\emph{Source and licence.} On the planar Taylor-Couette system and related exterior problems. Version arXiv:2406.14960v2 (CC BY 4.0), distributed by arXiv under the Creative Commons Attribution 4.0 International licence, http://creativecommons.org/licenses/by/4.0/, as recorded in the arXiv licence metadata for this exact version and shown on the abstract page https://arxiv.org/abs/2406.14960v2. Full licence notice: \texttt{licenses/arxiv-2406.14960-CC-BY-4.0.txt}.\par\smallskip

\noindent\textit{Editorial scope.} Counted unit: the article's theorem incompletepressure (source0.tex lines 653--660). The article's complete original proof of it is delivered with it (source0.tex lines 1025--1083, one \texttt{proof} environment of 59 lines, which the article places away from the statement). No mathematics is written, repaired or re-derived: the statement and the proof are the article's own lines, in the article's own notation, and no retained line is substituted or re-flowed.  Retained uncounted as the source-local context the proof needs: lines 135--142; lines 599--602; lines 736--739; lines 762--765; lines 1019--1021; lines 1296--1298; lines 1316--1318.  Nothing was omitted from any retained span, so the package carries no omission record.  No retained line contains a citation, so the wrapper carries no bibliography and no bibliography entry.  No retained line contains a figure environment, an image-inclusion command or drawing code, so no image is delivered and the package declares no asset dependency.  Editorial formatting only, in addition: an AMS \texttt{amsart} preamble that restates the article's own declarations and macros used by the retained lines, namely the environments \texttt{theorem} on separate counters, together with the standard packages those lines need.  The retained environments print in this document as Theorem 1 in order of appearance, whereas the article prints the same items with the numbering of its own full document.  The complete native source of the article as distributed in the arXiv e-print for this exact version is bundled unchanged as \texttt{sources/161123/source0.tex}.
\begin{equation}\label{nsstokes0}
\left\{
\begin{aligned}
& -\Delta u+(u\cdot\nabla)u+\nabla p=f \, ,\ \quad  \nabla\cdot u=0 \ \ \mbox{ in } \ \ \Omega_{R} \, , \\[3pt]
& u = \omega \, \widehat{\theta} \ \text{ on } \ \partial B_R \, , \qquad u = (0,0) \ \text{ on } \ \partial B_{1} \, .
\end{aligned}
\right.
\end{equation}

\begin{equation} \label{nstokesdebilstar}
\int_{\Omega_{R}} \nabla u \cdot \nabla \varphi + \int_{\Omega_{R}} (u \cdot \nabla) u \cdot \varphi  = \int_{\Omega_{R}} f \cdot \varphi \qquad
\forall \varphi \in \mathbb{H}_{\sigma}(\Omega_{R})\, .
\end{equation}

\begin{theorem} \label{nonuniquetheomain}
For any $\omega \geq 0$ and $f \in \mathcal{R}[L^{2}(\Omega_{R})]$, there exist infinitely many rotationally-invariant incomplete solutions to
\eqref{nsstokes0}.
\end{theorem}

\begin{equation}\label{lambdastarwf}
\int_{\Omega_R}\nabla V_{*}\cdot\nabla\varphi-\int_{\Omega_{R}}Q_{*}(\nabla\cdot\varphi)=0\qquad\forall\varphi\in
\mathbb{H}(\Omega_R)
\end{equation}

\begin{equation} \label{umbral1cor3}
\dfrac{\|f\|_{L^2(\Omega_{R})}}{\sqrt{\lambda(\Omega_{R})}} + K_{R}(\omega, \Phi) < \mathcal{S}_4  \, .
\end{equation}

\begin{equation} \label{sobolevconstants11}
\mathcal{S}_{p} \, \|v\|^{2}_{L^{p}(\Omega_{R})} \leq \|\nabla v\|^2_{L^2(\Omega_{R})} \qquad \forall v \in H_0^{1}(\Omega_{R}) \, .
\end{equation}

\begin{equation} \label{sobpoin}
\|v\|_{L^{2}(\Omega_{R})} \leq \dfrac{1}{\sqrt{\lambda(\Omega_{R})}} \|\nabla v\|_{L^2(\Omega_{R})} \qquad \forall v \in H_0^{1}(\Omega_{R}) \, .
\end{equation}

\begin{theorem} \label{incompletepressure}
For $\omega\ge0$ and $f \in L^2(\Omega_R)$, let $u \in H^{1}_{*,\sigma}(\Omega_{R})$ be an incomplete solution to \eqref{nsstokes0}.
Then, there exists a unique scalar pressure $p \in L^{2}(\Omega_{R})/\mathbb{R}$ such that
\begin{equation} \label{nstokesdebilstarpres}
\int_{\Omega_{R}} \nabla u \cdot \nabla \varphi + \int_{\Omega_{R}} (u \cdot \nabla) u \cdot \varphi  - \int_{\Omega_{R}} p (\nabla \cdot \varphi)
= \int_{\Omega_{R}} f \cdot \varphi \qquad \forall \varphi \in \mathbb{H}(\Omega_{R}) \, .
\end{equation}
\end{theorem}

\begin{proof} From the proof of Theorem
\ref{nonuniquetheomain} we know that there exists at least one rotationally-invariant incomplete solution $V_{\Phi} \in
\mathcal{R}[H^{1}_{*,\sigma}(\Omega_{R})]$ of \eqref{nsstokes0} having flux $\Phi$. Then, Theorem \ref{incompletepressure} ensures the existence
of a uniquely associated scalar pressure $Q_\Phi\in L^{2}(\Omega_{R})/\mathbb{R}$.\par
Concerning the first item of the statement, suppose that $V \in \mathcal{R}[H^{1}_{*,\sigma}(\Omega_{R})]$ is another incomplete solution to
\eqref{nsstokes0} having flux $\Phi$. Repeating the argument of the proof of Theorem \ref{nonuniquetheomain}, we deduce there exist two functions
$G_{\Phi}, G \in H^{1}(1,R)$ such that $G_{\Phi}(1)=G(1)=0$, $G_{\Phi}(R)=G(R)=\omega$ and
\begin{equation} \label{flux11}
V_{\Phi}(\xi) = G_{\Phi}(\rho) \widehat{\theta} \qquad \text{and} \qquad  V(\xi) = G(\rho) \widehat{\theta} \qquad \text{for a.e. }\xi \in
\Omega_{R} \, .
\end{equation}
Let $z := V - V_{\Phi} \in \mathbb{H}_{\sigma}(\Omega_{R})$ and subtract the equations \eqref{nstokesdebilstar} corresponding to $V$
and $V_{\Phi}$, thereby obtaining
$$
\int_{\Omega_{R}} \nabla z \cdot \nabla \varphi + \int_{\Omega_{R}} (V \cdot \nabla) z \cdot \varphi + \int_{\Omega_{R}} (z \cdot \nabla) V_{\Phi}
\cdot \varphi = 0 \qquad \forall \varphi \in \mathbb{H}_{\sigma}(\Omega_{R}) \, .
$$
By taking $\varphi=z$, integrating by parts, and using \eqref{flux11}, we get
$$
\int_{\Omega_{R}} (V \cdot \nabla) z \cdot z = 0 \qquad \text{and} \qquad \int_{\Omega_{R}} (z \cdot \nabla) V_{\Phi} \cdot z = -
\int_{\Omega_{R}} (z \cdot \nabla) z \cdot V_{\Phi} = 0 \, ,
$$
and we deduce that $z=0$ in $\Omega_{R}$. Then, Theorem \ref{incompletepressure} ensures that the associated pressures coincide (up to an additive constant)
a.e.\ in $\Omega_{R}$ as well.
\par
Concerning the second item of the statement, by taking $\varphi = V_{\Phi} - V_{*} \in \mathbb{H}_{\sigma}(\Omega_{R})$ in both weak formulations
\eqref{nstokesdebilstar} and \eqref{lambdastarwf}, and noticing that
$$
(V_{\Phi} \cdot \nabla) V_{\Phi} \cdot (V_{\Phi} - V_{*}) =0 \quad \text{a.e. in} \ \Omega_{R} \, ,
$$
we deduce
$$
\| \nabla V_{\Phi} - \nabla V_{*} \|^{2}_{L^{2}(\Omega_{R})} = \int_{\Omega_{R}} f \cdot (V_{\Phi} - V_{*}) \leq
\dfrac{\|f\|_{L^2(\Omega_{R})}}{\sqrt{\lambda(\Omega_{R})}} \| \nabla V_{\Phi} - \nabla V_{*} \|_{L^{2}(\Omega_{R})} \, ,
$$
as a consequence of H\"older's inequality and \eqref{sobpoin}. Therefore,
\begin{equation} \label{flux2}
\| \nabla V_{\Phi} \|_{L^{2}(\Omega_{R})} \leq \| \nabla V_{\Phi} - \nabla V_{*} \|_{L^{2}(\Omega_{R})} + \| \nabla V_{*} \|_{L^{2}(\Omega_{R})}
\leq \dfrac{\|f\|_{L^2(\Omega_{R})}}{\sqrt{\lambda(\Omega_{R})}} + \| \nabla V_{*} \|_{L^{2}(\Omega_{R})} \, .
\end{equation}
Now, suppose that $V \in H^{1}_{*,\sigma}(\Omega_{R})$ is another incomplete solution to  \eqref{nsstokes0} (not necessarily rotationally
invariant) having flux $\Phi$. Let $z := V - V_{\Phi} \in \mathbb{H}_{\sigma}(\Omega_{R})$ and subtract the equations
\eqref{nstokesdebilstar} corresponding to $V$ and $V_{\Phi}$, thereby obtaining
$$
\int_{\Omega_{R}} \nabla z \cdot \nabla \varphi + \int_{\Omega_{R}} (V \cdot \nabla) z \cdot \varphi + \int_{\Omega_{R}} (z \cdot \nabla) V_{\Phi}
\cdot \varphi = 0 \qquad \forall \varphi \in \mathbb{H}_{\sigma}(\Omega_{R}) \, .
$$
By taking $\varphi=z$ and noticing (after an integration by parts) that
$$
\int_{\Omega_{R}} (V \cdot \nabla) z \cdot z = 0 \, ,
$$
we deduce, from the H\"older inequality, from \eqref{sobolevconstants11} and \eqref{flux2}, that
$$
\|\nabla z\|^2_{L^2(\Omega_{R})} = - \int_{\Omega_{R}} (z \cdot \nabla)  V_{\Phi} \cdot z \leq \left(
\dfrac{\|f\|_{L^2(\Omega_{R})}}{\sqrt{\lambda(\Omega_{R})}} + \| \nabla V_{*} \|_{L^{2}(\Omega_{R})} \right) \frac{\|\nabla
z\|^2_{L^2(\Omega_{R})}}{\mathcal{S}_4}  \, ,
$$
so $z = 0$ a.e.\ in $\Omega_{R}$ whenever \eqref{umbral1cor3} holds.
\end{proof}

\end{document}
